\documentclass[11pt,a4paper]{article}

% ------------------------------------------------------------------
%  Kenya Informatics Olympiad -- Round 1
%  Set \answerstrue (or compile with \def\withanswers{}) to print
%  the marking scheme after the paper.
% ------------------------------------------------------------------
\newif\ifanswers
\answersfalse
\ifdefined\withanswers \answerstrue \fi

\usepackage[margin=2.2cm, headheight=14pt]{geometry}
\usepackage[T1]{fontenc}
\usepackage{lmodern}
\usepackage{amsmath,amssymb}
\usepackage{enumitem}
\usepackage{booktabs}
\usepackage{array}
\usepackage[dvipsnames]{xcolor}
\usepackage{fancyhdr}
\usepackage{lastpage}
\usepackage{titlesec}
\usepackage[most]{tcolorbox}
\usepackage[hidelinks]{hyperref}

% ---------- colours ----------
\definecolor{kioblack}{HTML}{1B1B1B}
\definecolor{kiored}{HTML}{B22222}
\definecolor{kiogreen}{HTML}{006B3F}
\definecolor{kiogrey}{HTML}{F2F2F2}

% ---------- header / footer ----------
\pagestyle{fancy}
\fancyhf{}
\fancyhead[L]{\small\textsc{Kenya Informatics Olympiad 2026}}
\fancyhead[R]{\small Round 1\ifanswers\ --- Marking Scheme\fi}
\fancyfoot[C]{\small Page \thepage\ of \pageref{LastPage}}
\renewcommand{\headrulewidth}{0.4pt}

% ---------- section style ----------
\titleformat{\section}
  {\large\bfseries\color{kiogreen}}{}{0pt}
  {}[\vspace{-0.6em}\color{kiogreen}\rule{\linewidth}{0.8pt}]
\titlespacing*{\section}{0pt}{1.6em}{0.9em}

% ---------- problems: numbered 1, 2, 3, ... with no titles ----------
\newcounter{prob}
\newenvironment{problem}[1]% #1 = marks
  {\refstepcounter{prob}\par\medskip
   \noindent{\bfseries\large \theprob.}\hfill\textit{[#1 marks]}\par
   \nopagebreak\smallskip}
  {\par\medskip}

\newcommand{\answerline}[1][Answer]{%
  \par\nopagebreak\smallskip\nopagebreak\noindent\textbf{#1:}\ \hrulefill\par}

\makeatletter
\newenvironment{parts}
  {\begin{enumerate}[label=(\alph*), leftmargin=2em, itemsep=0.35em]}
  {\@endparpenalty\@M\end{enumerate}}% keep the answer line with the parts
\makeatother

% ---------- multiple-choice options ----------
\newcommand{\mcq}[5]{%
  \par\nopagebreak\smallskip\noindent
  \begin{tabular*}{\linewidth}{@{\extracolsep{\fill}}lllll@{}}
  \textbf{(A)} #1 & \textbf{(B)} #2 & \textbf{(C)} #3 & \textbf{(D)} #4 & \textbf{(E)} #5
  \end{tabular*}\par}
\newenvironment{mcqlong}
  {\begin{enumerate}[label=\textbf{(\Alph*)}, leftmargin=2.4em, itemsep=0pt]}
  {\end{enumerate}}

% ---------- Section C boxes ----------
\newtcolorbox{example}{colback=white, colframe=kiogreen!60, boxrule=0.4pt,
  arc=2pt, left=5pt, right=5pt, top=3pt, bottom=3pt,
  title=\textbf{Example}, fonttitle=\small\bfseries, coltitle=white,
  colbacktitle=kiogreen!70}
\newcommand{\task}{\par\smallskip\noindent\textbf{Your task.} Describe an
  efficient algorithm for this problem. Explain the key idea, justify why it
  is correct, and state its time complexity. You do not need to write code.\par}

% ---------- solution environment ----------
\newtcolorbox{solution}[1][]{
  enhanced, breakable, colback=kiogrey, colframe=kiogreen,
  boxrule=0.5pt, arc=2pt, left=6pt, right=6pt, top=4pt, bottom=4pt,
  fonttitle=\bfseries, title={#1}}

\begin{document}

% ==================================================================
%  COVER
% ==================================================================
\begin{center}
  {\color{kioblack}\rule{\linewidth}{3pt}}\\[-2.2pt]
  {\color{kiored}\rule{\linewidth}{3pt}}\\[-2.2pt]
  {\color{kiogreen}\rule{\linewidth}{3pt}}\\[1.2em]
  {\LARGE\bfseries Kenya Informatics Olympiad 2026}\\[0.5em]
  {\Large Round 1 \ifanswers --- Marking Scheme\else --- Problem Paper\fi}\\[1em]
  {\color{kiogreen}\rule{\linewidth}{3pt}}\\[-2.2pt]
  {\color{kiored}\rule{\linewidth}{3pt}}\\[-2.2pt]
  {\color{kioblack}\rule{\linewidth}{3pt}}
\end{center}

\ifanswers\else
\vspace{0.6em}
\noindent
\begin{tabular}{@{}p{0.48\linewidth}p{0.48\linewidth}@{}}
Name: \hrulefill & School: \hrulefill \\[0.8em]
County: \hrulefill & Candidate No.: \hrulefill
\end{tabular}
\fi

\subsection*{Instructions}
\begin{enumerate}[leftmargin=2em, itemsep=0.2em]
  \item Time allowed: \textbf{3 hours}.
  \item The paper has \textbf{20 problems} in three sections, worth
        \textbf{100 marks} in total.
  \item \textbf{Section A} (Problems 1--7) is multiple choice. Each problem
        has exactly one correct option. Write its letter in the answer grid.
        Wrong answers do not lose marks.
  \item \textbf{Section B} (Problems 8--14) needs written answers. Show your
        working in the answer booklet; working earns partial credit.
  \item \textbf{Section C} (Problems 15--20) asks you to \emph{describe
        algorithms}. Write clear steps or pseudocode, explain why your method
        is correct, and give its running time. You do not need to write code
        in any particular programming language.
  \item Calculators and electronic devices are \textbf{not} allowed.
\end{enumerate}

\begin{tcolorbox}[colback=kiogrey, colframe=kioblack, boxrule=0.5pt,
                  arc=2pt, title=\textbf{Old British money (Sections A and B)},
                  fonttitle=\bfseries]
Several problems use pre-decimal British currency.
\begin{center}
\begin{tabular}{@{}ll@{\qquad}ll@{}}
\toprule
1 pound (\pounds1) & = 20 shillings (20s) & 1 crown & = 5s \\
1 shilling (1s)    & = 12 pence (12d)     & 1 half-crown & = 2s\,6d \\
1 penny (1d)       & = 4 farthings        & 1 florin & = 2s \\
1 half-sovereign   & = 10s                & 1 double florin & = 4s \\
\bottomrule
\end{tabular}
\end{center}
So, for example, 8s\,4d $= 8\times 12 + 4 = 100$ pence.
\end{tcolorbox}

% ==================================================================
\section*{Section A \quad Multiple Choice \hfill\normalsize (7 problems, 2 marks each)}
% ==================================================================

\ifanswers\else
\noindent\textbf{Answer grid.} Write one letter (A--E) in each box.
\begin{center}
\renewcommand{\arraystretch}{1.8}
\begin{tabular}{|c|*{7}{>{\centering\arraybackslash}p{1.1cm}|}}
\hline
\textbf{Problem} & 1 & 2 & 3 & 4 & 5 & 6 & 7 \\ \hline
\textbf{Answer}  &   &   &   &   &   &   &   \\ \hline
\end{tabular}
\end{center}
\fi

\begin{problem}{2}
A gentleman hands over a crown at the post office. He asks for some
twopenny stamps, six times as many penny stamps, and the rest of the
money in twopence-halfpenny stamps. How many stamps of each kind
(twopenny, penny, twopence-halfpenny) does he receive?
\mcq{4, 24, 11}{5, 30, 8}{3, 18, 14}{6, 36, 5}{2, 12, 17}
\end{problem}

\begin{problem}{2}
A gentleman meets three people asking for help, one after another. To the
first he gives 1d more than half the money in his pocket; to the second,
2d more than half of what is left; to the third, 3d more than half of what
is then left. He finishes with exactly 1d. How much did he start with?
\mcq{2s\,6d}{3s\,0d}{3s\,6d}{4s\,2d}{5s\,0d}
\end{problem}

\begin{problem}{2}
A group of cyclists agree to split a bill of \pounds4 equally. By the time
the bill is paid, two of them have left, so each cyclist who remains pays
2s more than his original share. How many cyclists were there at the start?
\mcq{5}{6}{8}{10}{12}
\end{problem}

\begin{problem}{2}
A man buys exactly 100 eggs of three qualities, costing 5d, 1d and
$\tfrac12$d each. He buys at least one egg of every quality, pays exactly
8s\,4d, and buys the same number of eggs of two of the qualities. How many
eggs does he buy at 5d, 1d and $\tfrac12$d respectively?
\mcq{10, 10, 80}{8, 20, 72}{12, 12, 76}{5, 5, 90}{20, 20, 60}
\end{problem}

\begin{problem}{2}
A grocer weighs out two 1-pound parcels of sugar per minute. A draper cuts
three 1-yard lengths of cloth per minute. Each must produce 48 items.
Customers interrupt them for 9 minutes in total between the two of them,
and the draper is interrupted for 17 times as long as the grocer. Which
statement is true?
\begin{mcqlong}
  \item The grocer finishes 1 minute before the draper.
  \item The grocer finishes 8 minutes before the draper.
  \item The draper finishes 1 minute before the grocer.
  \item The draper finishes $8\tfrac12$ minutes before the grocer.
  \item They finish at exactly the same time.
\end{mcqlong}
\end{problem}

\begin{problem}{2}
Hodge, Jakes and Durrant meet at a cattle market.
Hodge says to Jakes: ``I'll give you six of my pigs for one of your horses,
and then you'll have twice as many animals as I have.''
Durrant says to Hodge: ``I'll give you fourteen of my sheep for a horse,
and then you'll have three times as many animals as I.''
Jakes says to Durrant: ``I'll give you four cows for a horse, and then
you'll have six times as many animals as I have.''
Each offer is considered on its own, starting from the original herds.
How many animals did Hodge, Jakes and Durrant bring?
\mcq{7, 11, 21}{9, 3, 18}{11, 7, 21}{13, 11, 23}{21, 7, 11}
\end{problem}

\begin{problem}{2}
A man leaves instructions that exactly 55s be shared out once a year among
the poor of his parish: 18d to each woman and a half-crown to each man.
Every year at least one man and at least one woman must receive money, and
no two years may use the same numbers of men and women. For how many years
can the charity run?
\mcq{6}{7}{8}{11}{22}
\end{problem}

% ==================================================================
\section*{Section B \quad Written Answers \hfill\normalsize (7 problems, 36 marks)}
% ==================================================================

\begin{problem}{5}
Four brothers, John, William, Charles and Thomas, each have a money-box. If
John had 2s more, William 2s less, Charles twice as much and Thomas half as
much, all four would have the same amount.

Together the boxes hold 45s in exactly \textbf{six coins}. The coins
available are the sovereign (20s), half-sovereign (10s), crown (5s), double
florin (4s), half-crown (2s\,6d), florin (2s), shilling, sixpence and
threepence.
\begin{parts}
  \item How much money is in each box? \hfill [2]
  \item Which coins are in each box? \hfill [3]
\end{parts}
\answerline
\end{problem}

\begin{problem}{6}
Hiram B.\ Judkins has five droves of animals, all the same size, made up of
oxen, pigs and sheep (at least one of each kind overall). He sells every
animal to eight dealers, who each buy the same number of animals. He is paid
\$17 per ox, \$4 per pig and \$2 per sheep, and receives exactly \$301.
\begin{parts}
  \item What is the greatest number of animals he could have had? \hfill [2]
  \item Give \emph{every} possible combination of oxen, pigs and sheep for
        that total. \hfill [4]
\end{parts}
\answerline
\end{problem}

\begin{problem}{5}
An Englishman in New York buys goods costing 34 cents. All he has is a
dollar (100\,c), a three-cent piece and a two-cent piece. The shopkeeper has
only a half-dollar (50\,c) and a quarter-dollar (25\,c). Another customer
offers to help and puts down two dimes (10\,c each), a five-cent piece, a
two-cent piece and a one-cent piece.

The coins can be swapped around so that the Englishman pays exactly 34\,c
and the helpful customer ends up with exactly the money he started with.
\begin{parts}
  \item How much money does each of the three people hold once the deal is
        done? \hfill [1]
  \item Which coins does each person end up with? \hfill [4]
\end{parts}
\answerline
\end{problem}

\begin{problem}{5}
\begin{parts}
  \item A coin lies flat on a table. What is the greatest number of coins of
        the same size that can be placed flat around it so that each touches
        it (and none overlap)? Explain why. \hfill [1]
  \item A half-crown has diameter 32\,mm and a threepenny piece has diameter
        15\,mm. What is the greatest number of threepenny pieces that can lie
        flat on a half-crown, with none overlapping another or reaching
        beyond the edge? Justify your answer. \hfill [2]
  \item What is the largest diameter a threepenny piece could have for
        \emph{three} of them to fit on the half-crown? Give an exact answer.
        \hfill [2]
\end{parts}
\answerline
\end{problem}

\begin{problem}{5}
A millionaire wants to give away exactly \$1\,000\,000. Each gift must be
\$1 or a power of 7 (\$7, \$49, \$343, \dots), and no amount may be given
to more than six people.
\begin{parts}
  \item How many people receive a gift, and how many receive each amount?
        \hfill [2]
  \item Explain why there is only one way to do it. \hfill [2]
  \item How many people would receive a gift if he gave away
        \$1\,000\,000\,000 under the same rules? \hfill [1]
\end{parts}
\answerline
\end{problem}

\begin{problem}{5}
Take an amount \pounds$a$\;$s$s\;$d$d in which the number of pounds is
greater than the number of pence ($a > d$; zero shillings or pence are
allowed). Its \emph{reversal} is the amount with $d$ pounds, $s$ shillings
and $a$ pence, i.e.\ $240d + 12s + a$ pence (so $a$ may be 12 or more here).

Now apply this rule:
\begin{enumerate}[label=\arabic*., leftmargin=2em, itemsep=0pt]
  \item Subtract the reversal from the original amount, and write the
        difference in pounds, shillings and pence in the usual way
        (shillings 0--19, pence 0--11).
  \item Add the difference to its own reversal.
\end{enumerate}
For amounts with fewer than \pounds12 the result is always
\pounds12\;18s\;11d.
\begin{parts}
  \item Check the rule for \pounds7\;3s\;2d. \hfill [1]
  \item Show that the difference in step 1 is always $239(a-d)$ pence.
        \hfill [1]
  \item Without the ``fewer than \pounds12'' condition, what is the smallest
        amount for which the result is \emph{not} \pounds12\;18s\;11d?
        \hfill [1]
  \item What is the largest amount for which the result \emph{is}
        \pounds12\;18s\;11d? \hfill [2]
\end{parts}
\answerline
\end{problem}

\begin{problem}{5}
An excursion ticket costs 19s\,9d. The coins in use are: farthing,
halfpenny, penny, threepence, sixpence, shilling, florin, half-crown, double
florin, crown and half-sovereign. Two payments are different if they use a
different number of some coin (the order of the coins does not matter).
\begin{parts}
  \item In how many ways can you pay 3d? \hfill [1]
  \item In how many ways can you pay 6d? \hfill [2]
  \item Describe an efficient method for counting the ways to pay any amount,
        and explain why it does not count the same payment twice. \hfill [2]
\end{parts}
\answerline
\end{problem}

\clearpage
% ==================================================================
\section*{Section C \quad Algorithm Design \hfill\normalsize (6 problems, 50 marks)}
% ==================================================================

\noindent For each problem in this section, write an \textbf{algorithm
description}, not a program. A full answer explains the key observation,
lists the steps of the algorithm, argues why it gives the right answer, and
states the time complexity. The constraints tell you how fast your method
needs to be.

\begin{problem}{5}
You are given a string $s$ of $n$ lowercase letters and a letter $c$. In one
move you may pick any position $i$ and replace $s_i$ with the letter $c$.
(You may only ever write the letter $c$.)

Find the minimum number of moves needed to turn $s$ into a palindrome, i.e.\
a string that reads the same forwards and backwards.

\smallskip\noindent\textbf{Constraints.} $1 \le n \le 2\cdot 10^5$.
\begin{example}
$s = \texttt{kenya}$, $c = \texttt{a}$. The answer is $3$: change $s_1$, $s_2$
and $s_4$ to \texttt{a}, giving \texttt{aanaa}.
\end{example}
\task
\end{problem}

\begin{problem}{8}
The \emph{mode} of an array is the value that appears most often; if several
values tie, the mode is the largest of them. For example, the mode of
$[1,1,2]$ is $1$ and the mode of $[3,4]$ is $4$.

You are given an array $a$ of $n$ positive integers. You may rearrange its
elements in any order. Find an arrangement that makes the sum of the modes
of all $n$ prefixes as large as possible, and output that arrangement.

\smallskip\noindent\textbf{Constraints.} $1 \le n \le 2\cdot 10^5$,
$1 \le a_i \le 10^9$.
\begin{example}
$a = [2,5,2,5,5,1]$. The arrangement $[5,2,1,5,2,5]$ has mode $5$ for every
prefix, giving the maximum possible sum $30$.
\end{example}
\task
\end{problem}

\begin{problem}{9}
A treasury has $n$ piles of coins; pile $i$ holds $a_i$ coins. A pirate
holds a number $x$ and steals as follows:
\begin{enumerate}[leftmargin=2em, itemsep=0pt]
  \item He picks a pile $i$ with $a_i > 0$ and $\gcd(a_i, x) \ne 1$. If no
        such pile exists, he stops.
  \item Let $g = \gcd(a_i, x)$. He takes exactly $g$ coins from pile $i$
        (so $a_i$ decreases by $g$) and then sets $x = g$.
  \item He repeats from step 1.
\end{enumerate}
Find the maximum total number of coins he can steal.

\smallskip\noindent\textbf{Constraints.} $1 \le n \le 2\cdot 10^5$,
$1 \le a_i \le 3\cdot 10^5$, $2 \le x \le 3\cdot 10^5$.
\begin{example}
$x = 6$, $a = [4, 9, 10]$. The answer is $14$: he can empty the piles of
$4$ and $10$ coins.
\end{example}
\task
\end{problem}

\begin{problem}{9}
An array $b_1, \dots, b_m$ is a \emph{hill} if for some position $k$ it
strictly increases up to $b_k$ and strictly decreases after it:
$b_1 < \dots < b_k > b_{k+1} > \dots > b_m$. (A purely increasing or purely
decreasing array is also a hill.)

You are given an array $a$ of $n$ distinct integers. In one operation you
may swap $a_i$ and $a_{i+2}$ for any $i$. Decide whether $a$ can be turned
into a hill using any number of operations.

\smallskip\noindent\textbf{Constraints.} $1 \le n \le 2\cdot 10^5$.
\begin{example}
$[2,4,1,3,5]$: \textbf{YES} (for instance $[2,4,5,3,1]$ can be reached).
\quad $[1,3,2,4]$: \textbf{NO}.
\end{example}
\task
\end{problem}

\begin{problem}{9}
A binary string is \emph{clean} if all its characters are equal. In one move
you may choose a clean substring (a block of equal characters) and flip every
character in it ($0 \leftrightarrow 1$). The \emph{beauty} of a string is the
minimum number of moves needed to make the whole string clean. The
\emph{power} of a string is the sum of the beauties of all its substrings.

You are given a binary string $s$ of length $n$, followed by $q$ updates.
Each update gives a position $i$ and flips $s_i$. Output the power of $s$
before any update and after every update.

\smallskip\noindent\textbf{Constraints.} $1 \le n, q \le 2\cdot 10^5$.
\begin{example}
$s = \texttt{0110}$ has power $5$. After flipping position $1$ the string is
\texttt{1110}, whose power is $3$.
\end{example}
\task
\end{problem}

\begin{problem}{10}
For an array $b$ of $m$ integers, a \emph{transformation} does the following:
write down $b_i \oplus b_j$ for every pair $i < j$ (where $\oplus$ is bitwise
XOR), and replace the array with the $m$ smallest values written.

You are given an array $a$ of $n$ non-negative integers and $q$ queries.
Each query gives a number $k$; answer with $\max(a) - \min(a)$ after
applying the transformation $k$ times to the \emph{original} array.

\smallskip\noindent\textbf{Constraints.} $6 \le n \le 1000$,
$0 \le a_i < 2^{30}$, $1 \le q \le 2\cdot 10^5$, $0 \le k \le 10^9$.
\begin{example}
$a = [6,7,8,9,15,1]$. The arrays after $0,1,2,3$ transformations are
$[1,6,7,8,9,15]$, $[1,1,6,6,7,7]$, $[0,0,0,1,1,1]$ and $[0,0,0,0,0,0]$, so
the answers for $k = 0, 1, 2$ and any $k \ge 3$ are $14, 6, 1$ and $0$.
\end{example}
\task
\end{problem}

\begin{center}
\vspace{1em}
{\color{kiogreen}\rule{0.3\linewidth}{0.8pt}}\quad
\textsc{End of Paper}\quad
{\color{kiogreen}\rule{0.3\linewidth}{0.8pt}}
\end{center}

% ==================================================================
%  MARKING SCHEME (abridged here: only the Section A answer table is read by the converter;
%  keep your full marking scheme in this block as usual)
% ==================================================================
\ifanswers
\newpage
\section*{Marking Scheme}

\subsection*{Section A}
\begin{center}
\renewcommand{\arraystretch}{1.4}
\begin{tabular}{|c|*{7}{>{\centering\arraybackslash}p{1.1cm}|}}
\hline
\textbf{Problem} & 1 & 2 & 3 & 4 & 5 & 6 & 7 \\ \hline
\textbf{Answer}  & B & C & D & A & E & C & B \\ \hline
\end{tabular}
\end{center}
\fi

\end{document}
